\fancyhead[L]{\textbf{\textsf{66}}\quad \textbf{\textsf{\color{stewartcyan}CHƯƠNG 1}}\quad \textsf{HÀM SỐ VÀ MÔ HÌNH}}
\fancyhead[R]{}

\noindent
\begin{minipage}[t]{0.485\textwidth}
\small
\noindent\textbf{\textsf{\color{stewartcyan}55--56}}
\begin{enumerate}[label=(\alph*), itemjoin=\quad]
    \item Tìm tập xác định và tập giá trị của $f$?
    \item Giao điểm với trục hoành của đồ thị $f$ là gì?
    \item Vẽ đồ thị của $f$.
\end{enumerate}

\medskip
\noindent
\begin{tabularx}{\linewidth}{@{}X X@{}}
\textbf{55.} $f(x) = \ln x + 2$ & \textbf{56.} $f(x) = \ln(x - 1) - 1$
\end{tabularx}

\vspace{-1mm}
\noindent{\color{stewartcyan!60}\rule{\linewidth}{0.6pt}}
\vspace{1mm}

\noindent\textbf{\textsf{\color{stewartcyan}57--60}} Giải mỗi phương trình sau để tìm $x$. Hãy đưa ra cả giá trị chính xác và giá trị xấp xỉ thập phân, chính xác đến ba chữ số sau dấu phẩy.

\medskip
\noindent
\textbf{57.} (a) $\ln(4x + 2) = 3$ \hfill (b) $e^{2x-3} = 12$\hspace*{1.2cm}\null

\smallskip
\noindent
\textbf{58.} (a) $\log_2(x^2 - x - 1) = 2$ \hfill (b) $1 + e^{4x+1} = 20$\hspace*{0.7cm}\null

\smallskip
\noindent
\textbf{59.} (a) $\ln x + \ln(x - 1) = 0$ \hfill (b) $5^{1-2x} = 9$\hspace*{1.35cm}\null

\smallskip
\noindent
\textbf{60.} (a) $\ln(\ln x) = 0$ \hfill (b) $\dfrac{60}{1 + e^{-x}} = 4$\hspace*{1.1cm}\null

\vspace{0mm}
\noindent{\color{stewartcyan!60}\rule{\linewidth}{0.6pt}}
\vspace{1mm}

\noindent\textbf{\textsf{\color{stewartcyan}61--62}} Giải mỗi bất phương trình sau để tìm $x$.

\medskip
\noindent
\textbf{61.} (a) $\ln x < 0$ \hfill (b) $e^x > 5$\hspace*{2.2cm}\null

\smallskip
\noindent
\textbf{62.} (a) $1 < e^{3x-1} < 2$ \hfill (b) $1 - 2\ln x < 3$\hspace*{1.15cm}\null

\vspace{0mm}
\noindent{\color{stewartcyan!60}\rule{\linewidth}{0.6pt}}
\vspace{1mm}

\noindent\textbf{63.} (a) Tìm tập xác định của $f(x) = \ln(e^x - 3)$.\\
(b) Tìm $f^{-1}$ và tập xác định của nó.

\medskip
\noindent\textbf{64.} (a) Các giá trị của $e^{\ln 300}$ và $\ln(e^{300})$ là bao nhiêu?\\
(b) Sử dụng máy tính cầm tay của bạn để tính $e^{\ln 300}$ và $\ln(e^{300})$. Bạn nhận thấy điều gì? Bạn có thể giải thích tại sao máy tính lại gặp trục trặc không?

\medskip
\noindent\textbf{\color{stewartcyan}\fbox{\textbf{\textsf{T}}}}\;\textbf{65.} Vẽ đồ thị hàm số $f(x) = \sqrt{x^3 + x^2 + x + 1}$ và giải thích tại sao nó là hàm một-đối-một. Sau đó, hãy dùng một hệ thống đại số máy tính (CAS) để tìm một biểu thức tường minh cho $f^{-1}(x)$. (Hệ thống CAS sẽ đưa ra ba biểu thức khả dĩ. Hãy giải thích tại sao hai biểu thức trong số đó không phù hợp trong ngữ cảnh này.)

\medskip
\noindent\textbf{\color{stewartcyan}\fbox{\textbf{\textsf{T}}}}\;\textbf{66.} (a) Nếu $g(x) = x^6 + x^4$, $x \ge 0$, hãy dùng một hệ thống đại số máy tính để tìm một biểu thức cho $g^{-1}(x)$.\\
(b) Sử dụng biểu thức ở phần (a) để vẽ đồ thị của $y = g(x)$, $y = x$, và $y = g^{-1}(x)$ trên cùng một màn hình.

\medskip
\noindent\textbf{67.} Nếu một quần thể vi khuẩn ban đầu có 100 cá thể và tăng gấp đôi sau mỗi ba giờ, thì số lượng vi khuẩn sau $t$ giờ là $n = f(t) = 100 \cdot 2^{t/3}$.\\
(a) Tìm hàm ngược của hàm số này và giải thích ý nghĩa của nó.\\
(b) Khi nào thì quần thể sẽ đạt đến $50.000$ cá thể?

\medskip
\noindent\textbf{68.} Cơ sở Đánh lửa Quốc gia (NIF) tại Phòng thí nghiệm Quốc gia Lawrence Livermore duy trì cơ sở laser lớn nhất thế giới. Các tia laser, vốn được dùng để khởi động phản ứng nhiệt hạch, được cung cấp năng lượng bởi một hệ thống tụ điện lưu trữ tổng cộng khoảng 400 megajoule năng lượng. Khi các tia laser được
\end{minipage}%
\hfill
\begin{minipage}[t]{0.485\textwidth}
\small
\noindent
phát hỏa, các tụ điện sẽ phóng hết điện và ngay lập tức bắt đầu nạp lại. Điện tích $Q$ của các tụ điện sau $t$ giây kể từ khi phóng điện được cho bởi
\[
Q(t) = Q_0\left(1 - e^{-t/a}\right)
\]
(Điện tích nạp cực đại là $Q_0$ và $t$ được tính bằng giây.)\\
(a) Tìm công thức cho hàm ngược của hàm số này và giải thích ý nghĩa của nó.\\
(b) Mất bao lâu để nạp lại các tụ điện đạt đến 90\% dung lượng nếu $a = 50$?

\vspace{-1mm}
\noindent{\color{stewartcyan!60}\rule{\linewidth}{0.6pt}}
\vspace{1mm}

\noindent\textbf{\textsf{\color{stewartcyan}69--74}} Tìm giá trị chính xác của mỗi biểu thức sau.

\medskip
\noindent
\textbf{69.} (a) $\cos^{-1}(-1)$ \hfill (b) $\sin^{-1}(0{,}5)$\hspace*{1.3cm}\null

\smallskip
\noindent
\textbf{70.} (a) $\tan^{-1}\sqrt{3}$ \hfill (b) $\arctan(-1)$\hspace*{1.2cm}\null

\smallskip
\noindent
\textbf{71.} (a) $\csc^{-1}\sqrt{2}$ \hfill (b) $\arcsin 1$\hspace*{1.75cm}\null

\smallskip
\noindent
\textbf{72.} (a) $\sin^{-1}(-1/\sqrt{2})$ \hfill (b) $\cos^{-1}(\sqrt{3}/2)$\hspace*{1.1cm}\null

\smallskip
\noindent
\textbf{73.} (a) $\cot^{-1}(-\sqrt{3})$ \hfill (b) $\sec^{-1} 2$\hspace*{1.8cm}\null

\smallskip
\noindent
\textbf{74.} (a) $\arcsin\bigl(\sin(5\pi/4)\bigr)$ \hfill (b) $\cos\left(2\sin^{-1}\left(\frac{5}{13}\right)\right)$\hspace*{0.5cm}\null

\vspace{0mm}
\noindent{\color{stewartcyan!60}\rule{\linewidth}{0.6pt}}
\vspace{1mm}

\noindent\textbf{75.} Chứng minh rằng $\cos(\sin^{-1} x) = \sqrt{1 - x^2}$.

\vspace{-1mm}
\noindent{\color{stewartcyan!60}\rule{\linewidth}{0.6pt}}
\vspace{1mm}

\noindent\textbf{\textsf{\color{stewartcyan}76--78}} Rút gọn biểu thức.

\medskip
\noindent
\textbf{76.} $\tan(\sin^{-1} x)$ \hfill \textbf{77.} $\sin(\tan^{-1} x)$ \hfill \textbf{78.} $\sin(2\arccos x)$

\vspace{0mm}
\noindent{\color{stewartcyan!60}\rule{\linewidth}{0.6pt}}
\vspace{1mm}

\noindent\textbf{\color{stewartcyan}\fbox{\begin{tikzpicture}[scale=0.18, baseline=-1mm]
\draw[stewartcyan, line width=0.6pt] (0,0) rectangle (1.1,0.9);
\draw[stewartcyan, line width=0.5pt] (0.15,0.2) .. controls (0.45,0.7) and (0.65,0.2) .. (0.95,0.7);
\end{tikzpicture}}}\;\textbf{\textsf{\color{stewartcyan}79--80}} Vẽ đồ thị các hàm số đã cho trên cùng một màn hình. Những đồ thị này có mối liên hệ như thế nào?

\medskip
\noindent\textbf{79.} $y = \sin x,\ -\pi/2 \le x \le \pi/2;\quad y = \sin^{-1} x;\quad y = x$

\smallskip
\noindent\textbf{80.} $y = \tan x,\ -\pi/2 < x < \pi/2;\quad y = \tan^{-1} x;\quad y = x$

\vspace{0mm}
\noindent{\color{stewartcyan!60}\rule{\linewidth}{0.6pt}}
\vspace{1mm}

\noindent\textbf{81.} Tìm tập xác định và tập giá trị của hàm số
\[
g(x) = \sin^{-1}(3x + 1)
\]

\noindent\textbf{\color{stewartcyan}\fbox{\begin{tikzpicture}[scale=0.18, baseline=-1mm]
\draw[stewartcyan, line width=0.6pt] (0,0) rectangle (1.1,0.9);
\draw[stewartcyan, line width=0.5pt] (0.15,0.2) .. controls (0.45,0.7) and (0.65,0.2) .. (0.95,0.7);
\end{tikzpicture}}}\;\textbf{82.} (a) Vẽ đồ thị hàm số $f(x) = \sin(\sin^{-1} x)$ và giải thích hình dạng của đồ thị.\\
(b) Vẽ đồ thị hàm số $g(x) = \sin^{-1}(\sin x)$. Bạn giải thích thế nào về hình dáng của đồ thị này?

\medskip
\noindent\textbf{83.} (a) Nếu ta tịnh tiến một đường cong sang bên trái, điều gì sẽ xảy ra với hình chiếu đối xứng của nó qua đường thẳng $y = x$? Dưới góc nhìn hình học này, hãy tìm một biểu thức cho hàm ngược của $g(x) = f(x + c)$, trong đó $f$ là hàm một-đối-một.\\
(b) Tìm biểu thức cho hàm ngược của $h(x) = f(cx)$, trong đó $c \neq 0$.
\end{minipage}